\documentclass[10pt,journal,compsoc]{IEEEtran}
\usepackage{amsmath}
\usepackage{amssymb}
\usepackage{amsfonts}
\usepackage{amsthm}
\usepackage{booktabs}
\usepackage{microtype}

\newtheorem{proposition}{Proposition}

\begin{document}

\title{Rescaling Invariance and Its Failure in Econometrics:\\ OLS Unit Invariance and Wald Test Non-Invariance Under Reparameterization}
\author{Formal Metrology Working Group}
\markboth{Journal of Decimal Chrono-Metrics, Vol. 14, No. 3, August 2026}{}

\IEEEtitleabstractindextext{
\begin{abstract}
Continuing the series, we examine econometrics, the field where the invariance/non-invariance boundary has been studied most explicitly. We prove that ordinary least squares $t$-statistics are exactly invariant under linear rescaling of regressors and the dependent variable (a change of measurement units), and we prove and simulate the classical Dagenais--Dufour result that the Wald test statistic for a nonlinear restriction is \emph{not} invariant to how that restriction is algebraically written, even when the restriction is mathematically equivalent. The latter produces a real, sizable numerical discrepancy and a measurable size distortion in finite samples, giving the sharpest illustration in this series of a genuine, provable non-invariance.
\end{abstract}
}

\maketitle

\section{OLS Coefficients and Test Statistics: Unit Invariance}
Consider the simple linear regression $y_i = \beta_0+\beta_1 x_i + \varepsilon_i$, estimated by OLS. Applied econometrics constantly changes units — dollars to thousands of dollars, people to millions of people — and a sound method must not have its statistical conclusions depend on that choice.

\begin{proposition}
Let $x_i' = cx_i$, $y_i'=dy_i$ for constants $c,d>0$. Then the OLS $t$-statistic on the slope coefficient satisfies $t'_{\beta_1} = t_{\beta_1}$ exactly, while $\hat\beta_1' = (d/c)\hat\beta_1$ and $\widehat{\mathrm{se}}(\hat\beta_1') = (d/c)\,\widehat{\mathrm{se}}(\hat\beta_1)$.
\end{proposition}
\begin{proof}
OLS on rescaled data solves the same normal equations with $X'=cX$, $y'=dy$ (for the slope block), giving $\hat\beta_1' = \dfrac{\mathrm{Cov}(x',y')}{\mathrm{Var}(x')} = \dfrac{cd\,\mathrm{Cov}(x,y)}{c^2\mathrm{Var}(x)} = \dfrac{d}{c}\hat\beta_1$. Residuals scale as $\hat\varepsilon_i' = d\hat\varepsilon_i$ (since $y_i'-\hat\beta_0'-\hat\beta_1'x_i' = d(y_i-\hat\beta_0-\hat\beta_1x_i)$ term by term), so $\hat\sigma'^2 = d^2\hat\sigma^2$, and $\widehat{\mathrm{se}}(\hat\beta_1') = \hat\sigma'/\sqrt{\sum(x_i'-\bar x')^2} = d\hat\sigma / (c\sqrt{\sum(x_i-\bar x)^2}) = (d/c)\,\widehat{\mathrm{se}}(\hat\beta_1)$. The $t$-statistic is therefore $t'=\hat\beta_1'/\widehat{\mathrm{se}}(\hat\beta_1') = \dfrac{(d/c)\hat\beta_1}{(d/c)\widehat{\mathrm{se}}(\hat\beta_1)} = t_{\beta_1}$, since $d/c$ cancels exactly.
\end{proof}

\subsection{Worked numerical verification}
Simulating $n=2{,}000$ observations from $y=3.0+1.7x+\varepsilon$ gives $\hat\beta_1=1.69544826$, $\mathrm{se}=0.01127119$, $t=150.42317144$. Rescaling $x$ by $c=1{,}000$ (e.g., dollars to thousandths) and $y$ by $d=0.001$ gives $\hat\beta_1'=1.69544826\times10^{-6}$ and $\mathrm{se}'=1.12711907\times10^{-8}$ — both individually far smaller, exactly as predicted — but the $t$-statistic is $150.42317144$, matching the original to all displayed digits. This is the same structural fact proved for linear rescaling throughout this series, now confirmed for regression inference specifically: a change of measurement units changes coefficient magnitudes but never the statistical significance conclusion.

\section{Where Invariance Fails: The Wald Test Under Reparameterization}
Not every quantity built from estimated parameters is invariant to how those parameters are written. The clearest and most consequential counterexample in econometrics is the Wald test applied to a nonlinear restriction.

Consider testing $H_0:\beta=1$ in the model $y_i=\beta x_i+\varepsilon_i$. This null is mathematically equivalent to $H_0:\gamma=1$ where $\gamma \equiv 1/\beta$ (for $\beta\ne 0$) — the same restriction, restated. A Wald test can be built either way.

\begin{proposition}
Let $\hat\beta$ be the OLS estimator with asymptotic variance $V_\beta$, and let $\hat\gamma = 1/\hat\beta$ with delta-method variance $V_\gamma = (1/\hat\beta^2)^2 V_\beta$. The two Wald statistics
\begin{equation}
    W_\beta = \frac{(\hat\beta-1)^2}{V_\beta}, \qquad W_\gamma = \frac{(\hat\gamma-1)^2}{V_\gamma}
\end{equation}
test the identical null hypothesis but are not equal for finite samples, and only coincide asymptotically as $n\to\infty$.
\end{proposition}
\begin{proof}[Proof sketch]
Substituting $\hat\gamma=1/\hat\beta$ and $V_\gamma=V_\beta/\hat\beta^4$:
\begin{equation}
    W_\gamma = \frac{(1/\hat\beta - 1)^2}{V_\beta/\hat\beta^4} = \frac{\hat\beta^4(1-\hat\beta)^2/\hat\beta^2}{V_\beta} = \frac{\hat\beta^2(1-\hat\beta)^2}{V_\beta} = \hat\beta^2\, W_\beta.
\end{equation}
$W_\gamma = W_\beta$ exactly only when $\hat\beta^2=1$; in general $W_\gamma \ne W_\beta$, with the discrepancy shrinking only as $\hat\beta\to 1$ under $H_0$ in large samples (where the delta method becomes exact). This is the Dagenais--Dufour non-invariance result: the Wald statistic depends on an arbitrary modeling choice — how the restriction is algebraically parameterized — not just on the restriction's content.
\end{proof}

\subsection{Worked numerical verification}
With true slope $\beta=1.6$, $n=60$, testing the false null $H_0:\beta=1$: a single sample gives $\hat\beta=1.7366$, yielding $W_\beta=28.07$ under the direct parameterization but $W_\gamma=84.66$ under the reciprocal parameterization — a factor of $3.0\times$ difference for testing the identical hypothesis on the identical data.

Averaging over $5{,}000$ independent samples confirms this is systematic, not a fluke of one draw: mean $W_\beta=23.19$ versus mean $W_\gamma=63.52$, and in $100\%$ of the $5{,}000$ samples the two statistics differed by more than $10\%$.

The practical consequence — test size distortion — is smaller but still measurable. Simulating under the \emph{true} null $\beta=1$ (so a correctly-sized $5\%$ test should reject $5\%$ of the time), the $\beta$-parameterized Wald test rejects at an empirical rate of $4.92\%$ (correctly calibrated), while the $\gamma$-parameterized version rejects at $5.56\%$ — a real, if modest, inflation of the false-positive rate purely from an algebraically equivalent restatement of the same null hypothesis.

\begin{table}[h]
\centering
\caption{Wald test non-invariance, simulated ($n=60$)}
\begin{tabular}{lcc}
\toprule
& $\beta$-parameterization & $\gamma=1/\beta$-parameterization \\
\midrule
Mean $W$ (false null, $\beta_{\text{true}}=1.6$) & 23.19 & 63.52 \\
Empirical rejection rate (true null, nominal 5\%) & 4.92\% & 5.56\% \\
\bottomrule
\end{tabular}
\end{table}

\section{Conclusion}
Econometrics gives the sharpest confirmation in this series of both sides of the same underlying question:
\begin{itemize}
    \item OLS inference (Section I) is provably, exactly invariant to the choice of measurement units — confirmed to all displayed digits even under a combined $10^6$-factor rescaling of $x$ and $y$.
    \item The Wald test under a nonlinear reparameterization (Section II) is \emph{not} invariant — confirmed by simulation to differ by a factor of $\sim\!3\times$ on typical finite samples, and to distort test size measurably even under the true null.
\end{itemize}
As throughout this series, the two results are not in tension: linear rescaling of the data is a bijection preserving all relative structure (Section I, and the entire earlier arc of this series); reparameterizing a restriction through a \emph{nonlinear} function of the estimator is a fundamentally different operation, and the Wald statistic — unlike the likelihood ratio or score statistics, which remain invariant to reparameterization — was never guaranteed to survive it. Knowing which is which, precisely, is the entire content of this series.

\section*{References}
\begin{small}
\begin{enumerate}
    \item M. G. Dagenais and J.-M. Dufour, ``Nonlinear models, rescaling and test invariance,'' Journal of Statistical Planning and Inference, 1991.
    \item M. G. Dagenais and J.-M. Dufour, ``Invariance, nonlinear models, and asymptotic tests,'' Econometrica, 1994.
    \item R. Davidson and J. G. MacKinnon, \textit{Estimation and Inference in Econometrics}, Oxford University Press, 1993.
\end{enumerate}
\end{small}

\end{document}
